% !TeX root = Master.tex
Monte Carlo simulations are widely used to evaluate causal estimators and guide estimator selection in applied research. In empirical Monte Carlo studies, simulation models are often fitted to observed data and used to generate benchmark datasets with a specified target value for the causal estimand. However, it may be unclear whether the resulting benchmarks reproduce the conclusions that would have been obtained under the original data-generating process. In this paper, we formalize benchmarks in empirical Monte Carlo studies as pairs consisting of a dataset law and a reference estimand value, which are designed to approximate an unavailable target benchmark under the unknown data-generating process. To illustrate our framework and gain empirical evidence, we employ two established collections of data-generating designs from the causal-inference literature to compare outcome simulators in their ability to preserve estimator rankings and guide estimator selection. We focus particularly on bias arising from sharing model families with candidate estimators. We assess agreement between target and simulator-induced benchmarks using rank correlations. Results show that simulators often preserve rankings well, but with substantial variability across models and settings. Sharing model families with candidate estimators can distort their rankings, particularly when the simulation model struggles to capture the underlying outcome structure. In these cases, simulator-guided selection can improve on direct estimation after excluding aligned estimators, whereas direct estimation can be preferable when those models perform well. Finally, distances between observed and simulated outcome distributions were somewhat informative, particularly for ruling out very poorly performing simulators.
